\section{卷二 42-72}

\subsection{第42篇(可拉後裔)： 因“神在哪里”而引发的祷告 42:1-11}
\subsubsection{因渴想神，神殿和人不住问“神在哪里”，使我心憂悶煩躁 1-5}
\begin{table}[H]
\centering
\begin{tabular}{p{13.5cm}|p{3.700cm}}\hline\hline
\multicolumn{1}{c|}{憂悶的原因}&\multicolumn{1}{c}{憂悶煩躁應當仰望}\\\hline
神啊，我的心切慕你，如鹿切慕溪水。我的心渴想神，就是永生神，我幾時得朝見神呢？
&\multirow{3}{3.9cm}{我的心哪，你為何憂悶？為何在我裡面煩躁？應當仰望神，因他笑臉幫助我，我還要稱讚他。}\\\cline{1-1}
3 我晝夜以眼淚當飲食，人不住地對我說：「你的神在哪裡呢？」&\\\cline{1-1}
我從前與眾人同往，用歡呼稱讚的聲音，領他們到神的殿裡，大家守節。我追想這些事，我的心極其悲傷。&\\\hline
\end{tabular}
\end{table} 

\subsubsection{憂悶煩躁就要仰望稱讚 6-11}
6 我的神啊，我的心在我裡面憂悶，所以我從約旦地，從黑門嶺，從米薩山，記念你。
7 你的瀑布發聲，深淵就與深淵響應，你的波浪洪濤漫過我身。
8 白晝，耶和華必向我施慈愛；黑夜，我要歌頌、禱告賜我生命的神。
9 我要對神我的磐石說：「你為何忘記我呢？我為何因仇敵的欺壓時常哀痛呢？」
我的敵人辱罵我，好像打碎我的骨頭，不住地對我說：「你的神在哪裡呢？」
11我的心哪，你為何憂悶？為何在我裡面煩躁？應當仰望神，因我還要稱讚他，他是我臉上的光榮（幫助），是我的神。

\subsection{第43篇： 憂悶煩躁时的祷告 43:1-5}
\subsubsection{向神的祈求 1-3}
神啊，求你申我的冤，向不虔誠的國為我辨屈，求你救我脫離詭詐不義的人！
2 因為你是賜我力量的神。為何丟棄我呢？我為何因仇敵的欺壓時常哀痛呢？
3 求你發出你的亮光和真實，好引導我，帶我到你的聖山，到你的居所。
\subsubsection{我要做的回应 4-5}
我就走到神的祭壇，到我最喜樂的神那裡。神啊，我的神，我要彈琴稱讚你。
5 我的心哪，你為何憂悶？為何在我裡面煩躁？應當仰望神，因我還要稱讚他，他是我臉上的光榮（幫助），是我的神。

\subsection{第44篇(可拉後裔)： 提醒神古时的恩惠，并向被掳的百姓施恩 44:1-26}
\subsubsection{追念神古昔的作为 1-8}
\begin{table}[H]
\centering
\begin{tabular}{p{9.25cm}|p{8.50cm}}\hline\hline
\multicolumn{1}{c|}{神在古時我們列祖的日子所行的事1-3}&\multicolumn{1}{c}{我們終日因神誇耀永遠稱謝的事4-8}\\\hline
神啊，你在古時我們列祖的日子所行的事，我們親耳聽見了，我們的列祖也給我們述說過。
2 你曾用手趕出外邦人，卻栽培了我們列祖；你苦待列邦，卻叫我們列祖發達。
3 因為他們不是靠自己的刀劍得地土，也不是靠自己的膀臂得勝，乃是靠你的右手、你的膀臂和你臉上的亮光，因為你喜悅他們。
&神啊，你是我的王，求你出令使雅各得勝。
5 我們靠你要推倒我們的敵人，靠你的名要踐踏那起來攻擊我們的人。
6 因為我必不靠我的弓，我的刀也不能使我得勝。
7 唯你救了我們脫離敵人，使恨我們的人羞愧。
8 我們終日因神誇耀，還要永遠稱謝你的名。（細拉）\\\hline
\end{tabular}
\end{table} 

\subsubsection{陳述己身今時之苦，并救贖求神 9-26}
\begin{table}[H]
\centering
\begin{tabular}{p{7.5cm}|p{4.25cm}|p{4.500cm}}\hline\hline
\multicolumn{1}{c|}{述说神如何抛弃百姓 9-16}&\multicolumn{1}{c|}{困境中不忘記神17-21}&\multicolumn{1}{c}{求神帮助的祷告 22-26}\\\hline
但如今你丟棄了我們，使我們受辱，不和我們的軍兵同去。
你使我們向敵人轉身退後，那恨我們的人任意搶奪。
你使我們當做快要被吃的羊，把我們分散在列邦中。
你賣了你的子民也不賺利，所得的價值並不加添你的資財。
你使我們受鄰國的羞辱，被四圍的人嗤笑譏刺。
你使我們在列邦中做了笑談，使眾民向我們搖頭。
我的凌辱終日在我面前，我臉上的羞愧將我遮蔽，
都因那辱罵毀謗人的聲音，又因仇敵和報仇人的緣故。
&這都臨到我們身上，我們卻沒有忘記你，也沒有違背你的約。
我們的心沒有退後，我們的腳也沒有偏離你的路。
你在野狗之處壓傷我們，用死蔭遮蔽我們。
倘若我們忘了神的名，或向別神舉手，神豈不鑒察這事嗎？因為他曉得人心裡的隱祕。
&我們為你的緣故終日被殺，人看我們如將宰的羊。
主啊，求你睡醒！為何儘睡呢？求你興起，不要永遠丟棄我們！
你為何掩面，不顧我們所遭的苦難和所受的欺壓？
我們的性命伏於塵土，我們的肚腹緊貼地面。
求你起來幫助我們，憑你的慈愛救贖我們！\\\hline
\end{tabular}
\end{table} 



\subsection{第45篇(基督的婚宴-诗歌中的雅歌。可拉後裔)： 王与新妇快乐的婚宴 45:1-17}
\subsubsection{赞美王的荣美 1-8}
\begin{table}[H]
\centering
\begin{tabular}{p{2.25cm}|p{15.cm}}\hline\hline
\multicolumn{1}{c|}{用美辭讚美王}&\multicolumn{1}{c}{述说王榮威美德2-8}\\\hline
我心裡湧出美辭，我論到我為王做的事，我的舌頭是快手筆。
&你比世人更美，在你嘴裡滿有恩惠，所以神賜福給你，直到永遠。
3 大能者啊，願你腰間佩刀，大有榮耀和威嚴。
4 為真理、謙卑、公義赫然坐車前往，無不得勝，你的右手必顯明可畏的事。
5 你的箭鋒快，射中王敵之心，萬民仆倒在你以下。
6 神啊，你的寶座是永永遠遠的，你的國權是正直的。
7 你喜愛公義，恨惡罪惡，所以神，就是你的神，用喜樂油膏你，勝過膏你的同伴。
8 你的衣服都有沒藥、沉香、肉桂的香氣，象牙宮中有絲弦樂器的聲音使你歡喜。\\\hline
\end{tabular}
\end{table} 

\subsubsection{讚美王的新妇 9-17}
\begin{table}[H]
\centering
\begin{tabular}{p{7.75cm}|p{5.25cm}|p{3.5cm}}\hline\hline
\multicolumn{1}{c|}{新妇的尊貴9-12}&\multicolumn{1}{c}{新妇的穿着13-15}&\multicolumn{1}{c}{新妇的福气16-17}\\\hline
有君王的女兒在你尊貴婦女之中，王后佩戴俄斐金飾站在你右邊。
女子啊，你要聽，要想，要側耳而聽！不要記念你的民和你的父家，
王就羨慕你的美貌。因為他是你的主，你當敬拜他。
推羅的民（女子）必來送禮，民中的富足人也必向你求恩。
&王女在宮裡極其榮華，她的衣服是用金線繡的。
14 她要穿錦繡的衣服，被引到王前，隨從她的陪伴童女也要被帶到你面前。
15 她們要歡喜快樂被引導，她們要進入王宮。
&你的子孫要接續你的列祖，你要立他們在全地做王。
我必叫你的名被萬代記念，所以萬民要永永遠遠稱謝你。\\\hline
\end{tabular}
\end{table} 


\subsection{第46篇(可拉後裔)： 萬軍之耶和華的同在是我們的避難所 46:1-11}
\subsubsection{神是我們的避難所 1-3}
\begin{table}[H]
\centering
\begin{tabular}{p{7.cm}|p{10.25cm}}\hline\hline
\multicolumn{1}{c|}{神是1}&\multicolumn{1}{c}{我們不害怕2-3}\\\hline
神是我們的避難所，是我們的力量，是我們在患難中隨時的幫助。
&所以地雖改變，山雖搖動到海心，
3 其中的水雖砰訇翻騰，山雖因海漲而戰抖，我們也不害怕。（細拉）\\\hline
\end{tabular}
\end{table} 

\subsubsection{神的城 4-7}
有一道河，這河的分汊使神的城歡喜，這城就是至高者居住的聖所。
神在其中，城必不動搖，到天一亮，神必幫助這城。
外邦喧嚷，列國動搖，神發聲，地便熔化。
萬軍之耶和華與我們同在，雅各的神是我們的避難所。（細拉）
\subsubsection{神的作为 8-11}
你們來看耶和華的作為，看他使地怎樣荒涼。
9 他止息刀兵，直到地極；他折弓斷槍，把戰車焚燒在火中。
10 你們要休息，要知道我是神！我必在外邦中被尊崇，在遍地上也被尊崇。
11 萬軍之耶和華與我們同在，雅各的神是我們的避難所。

\subsection{第47篇(可拉後裔)： 耶和華為全地之君，萬民都要向神歌頌 47:1-9}
\subsubsection{呼吁万民敬拜神 1-2}
\begin{itemize}
\itemsep-.45em  
\item 萬民哪，你們都要拍掌，要用誇勝的聲音向神呼喊
\item 因為耶和華至高者是可畏的，他是治理全地的大君王
\end{itemize}
\subsubsection{神在以色列的作为 3-4}
\begin{itemize}
\itemsep-.45em  
\item 他叫萬民服在我們以下
\item 又叫列邦服在我們腳下
\item 他為我們選擇產業就是他所愛之雅各的榮耀
\end{itemize}
\subsubsection{指教百姓如何敬拜神 5-9}
\begin{itemize}
\itemsep-.45em  
\item 神上升，有喊聲相送，有角聲相送
\item 你們要向神歌頌，歌頌！向我們王歌頌，歌頌
\item 你們要用悟性歌頌，因為神是全地的王
\item 列邦的君王聚集，要做亞伯拉罕之神的民，因为
\begin{itemize}
\itemsep-.45em  
\item 因為世界的盾牌是屬神的，他為至高
\item 神做王治理萬國，神坐在他的聖寶座上
\end{itemize}
\end{itemize}

\subsection{第48篇(可拉後裔)： 頌讚錫安之榮美 1-14}
\subsubsection{眾王見城喪膽逃跑 1-7}
耶和華本為大，在我們神的城中，在他的聖山上，該受大讚美！
2 錫安山，大君王的城，在北面居高華美，為全地所喜悅。
3 神在其宮中自顯為避難所。

4 看哪，眾王會合，一同經過。
5 他們見了這城，就驚奇喪膽，急忙逃跑。
6 他們在那裡被戰兢疼痛抓住，好像產難的婦人一樣。
7 神啊，你用東風打破他施的船隻。
\subsubsection{聖民在城中讚美享慈愛 8-14}
8 我們在萬軍之耶和華的城中，就是我們神的城中，所看見的正如我們所聽見的：神必堅立這城，直到永遠。（細拉）
9 神啊，我們在你的殿中想念你的慈愛。
10 神啊，你受的讚美正與你的名相稱，直到地極，你的右手滿了公義。
11 因你的判斷，錫安山應當歡喜，猶大的城邑(女子)應當快樂。
12 你們當周遊錫安，四圍旋繞，數點城樓，
13 細看她的外郭，察看她的宮殿，為要傳說到後代。
14 因為這神永永遠遠為我們的神，他必做我們引路的，直到死時。

\subsection{第49篇(可拉後裔)： 嘆世人恃財之愚,見人財室增榮不要懼怕 1-20}
\subsubsection{世人恃財之愚 1-11}
\begin{table}[H]
\centering
\begin{tabular}{p{7.25cm}|p{10.25cm}}\hline\hline
\multicolumn{1}{c|}{萬民都當聽的話1-5}&\multicolumn{1}{c}{世人恃財之愚6-11}\\\hline
萬民哪，你們都當聽這話。世上一切的居民，
2 無論上流下流，富足貧窮，都當留心聽。
3 我口要說智慧的言語，我心要想通達的道理。
4 我要側耳聽比喻，用琴解謎語。在患難的日子，奸惡隨我腳跟四面環繞我，我何必懼怕？&
6 那些倚仗財貨，自誇錢財多的人，
7 一個也無法贖自己的弟兄，也不能替他將贖價給神，
8 叫他長遠活著，不見朽壞，因為贖他生命的價值極貴，只可永遠罷休。 9 
10 他必見智慧人死，又見愚頑人和畜類人一同滅亡，將他們的財貨留給別人。
11 他們心裡思想他們的家室必永存，住宅必留到萬代，他們以自己的名稱自己的地。\\\hline
\end{tabular}
\end{table} 

\subsubsection{尊貴中不醒悟如死亡的畜類 12-20}
\begin{table}[H]
\centering
\begin{tabular}{p{7.cm}|p{10.5cm}}\hline\hline
\multicolumn{1}{c|}{後人愚昧的佩服12-14}&\multicolumn{1}{c}{見人財室增榮不要懼怕15-20}\\\hline
12 但人居尊貴中不能長久，如同死亡的畜類一樣。
13 他們行的這道本為自己的愚昧，但他們以後的人還佩服他們的話語！（細拉）
14 他們如同羊群派定下陰間，死亡必做他們的牧者。到了早晨，正直人必管轄他們。他們的美容必被陰間所滅，以致無處可存。
&15 只是神必救贖我的靈魂脫離陰間的權柄，因他必收納我。（細拉）
16 見人發財家室增榮的時候，你不要懼怕，
17 因為他死的時候，什麼也不能帶去，他的榮耀不能隨他下去。
18 他活著的時候，雖然自誇為有福（你若利己，人必誇獎你），
19 他仍必歸到他歷代的祖宗那裡，永不見光。
20 人在尊貴中而不醒悟，就如死亡的畜類一樣。\\\hline
\end{tabular}
\end{table} 



\subsection{第50篇(亞薩)： 神審判聖民惡人 50:1-23}
\subsubsection{神發言招呼天下聖民聚集為要審判 1-6}
大能者神耶和華已經發言招呼天下，從日出之地到日落之處。
2 從全美的錫安中，神已經發光了。
3 我們的神要來，決不閉口。有烈火在他面前吞滅，有暴風在他四圍大颳。
4 他招呼上天下地，為要審判他的民，
5 說：「招聚我的聖民到我這裡來，就是那些用祭物與我立約的人。」
6 諸天必表明他的公義，因為神是施行審判的。（細拉）
\subsubsection{警教選民 7-15}
\begin{table}[H]
\centering
\begin{tabular}{p{13.cm}|p{3.4cm}}\hline\hline
\multicolumn{1}{c|}{神不要的}&\multicolumn{1}{c}{神所要的}\\\hline
我的民哪，你們當聽我的話；以色列啊，我要勸誡你。我是神，是你的神！我並不因你的祭物責備你，你的燔祭常在我面前。我不從你家中取公牛，也不從你圈內取山羊。
&\multirow{2}{3.5cm}{你們要以\textcolor{red}{感謝}為祭獻於神，又要向至高者\textcolor{red}{還你的願}；並要在患難之日\textcolor{red}{求告}我，我必搭救你，你也要\textcolor{red}{榮耀}我。}\\\cline{1-1}
10 因為樹林中的百獸是我的，千山上的牲畜也是我的。
11 山中的飛鳥我都知道，野地的走獸也都屬我。
12 我若是飢餓，我不用告訴你，因為世界和其中所充滿的都是我的。
13 我豈吃公牛的肉呢？我豈喝山羊的血呢？&\\\hline
\end{tabular}
\end{table} 

\subsubsection{嚴責惡人 17-23}
\begin{table}[H]
\centering
\begin{tabular}{p{8.75cm}|p{8.00cm}}\hline\hline
\multicolumn{1}{c|}{惡人所作所說}&\multicolumn{1}{c}{神的審判}\\\hline
但神對惡人說：「你怎敢傳說我的律例，口中提到我的約呢？
17 其實你恨惡管教，將我的言語丟在背後。
18 你見了盜賊，就樂意與他同夥，又與行姦淫的人一同有份。
&\multirow{2}{8.1cm}{21 你行了這些事，我還閉口不言，你想我恰和你一樣。其實我要責備你，將這些事擺在你眼前。
22 你們忘記神的，要思想這事，免得我把你們撕碎，無人搭救。
23 凡以感謝獻上為祭的，便是榮耀我；那按正路而行的，我必使他得著我的救恩。」}\\\cline{1-1}
19 你口任說惡言，你舌編造詭詐。
20 你坐著毀謗你的兄弟，讒毀你親母的兒子。&\\\hline
\end{tabular}
\end{table} 

\subsection{第51篇(大衛)： 誠懇后悔流人血的罪，切求恩免 1-19}
\subsubsection{陈明自己是罪人，求神按豐盛的慈悲塗抹過犯 1-5}
大衛與拔示巴同室以後，先知拿單來見他。他作這詩，交於伶長。
1 神啊，\textcolor{red}{求}你按你的慈愛憐恤我，按你豐盛的慈悲塗抹我的過犯！
2 \textcolor{red}{求}你將我的罪孽洗除淨盡，並潔除我的罪！
3 因為我知道我的過犯，我的罪常在我面前。
4 我向你犯罪，唯獨得罪了你，在你眼前行了這惡，以致你責備我的時候顯為公義，判斷我的時候顯為清正。
5 我是在罪孽裡生的，在我母親懷胎的時候就有了罪。

\subsubsection{神喜愛內裡誠實，求神救我脫離流人血的罪 6-15}
6 你所喜愛的是內裡誠實，你在我隱密處必使我得智慧。
7 \textcolor{red}{求}你用牛膝草潔淨我，我就乾淨；求你洗滌我，我就比雪更白。
8 \textcolor{red}{求}你使我得聽歡喜快樂的聲音，使你所壓傷的骨頭可以踴躍。
9 \textcolor{red}{求}你掩面不看我的罪，塗抹我一切的罪孽。
10 神啊，\textcolor{red}{求}你為我造清潔的心，使我裡面重新有正直(堅定)的靈。
11 不要丟棄我，使我離開你的面，不要從我收回你的聖靈。
12 \textcolor{red}{求}你使我仍得救恩之樂，賜我樂意的靈扶持我。
13 我就把你的道指教有過犯的人，罪人必歸順你。
14 神啊，你是拯救我的神，\textcolor{red}{求}你救我脫離流人血的罪，我的舌頭就高聲歌唱你的公義。
15 主啊，\textcolor{red}{求}你使我嘴唇張開，我的口便傳揚讚美你的話。

\subsubsection{神喜愛的公義祭物是憂傷痛悔的心，求神善待錫安 16-19}
16 你本不喜愛祭物，若喜愛，我就獻上，燔祭你也不喜悅。
17 神所要的祭就是憂傷的靈，神啊，憂傷痛悔的心，你必不輕看。
18 \textcolor{red}{求}你隨你的美意善待錫安，建造耶路撒冷的城牆。
19 那時，你必喜愛公義的祭和燔祭，並全牲的燔祭。那時，人必將公牛獻在你壇上。

\subsection{第52篇(大衛)：嘆惡人（以東人多益）狂妄自誇 52:1-9}
\begin{table}[H]
\centering
\begin{tabular}{p{9.cm}|p{7.500cm}}\hline\hline
\multicolumn{1}{c|}{外人攻擊，求神申冤}&\multicolumn{1}{c}{我要獻祭稱讚神}\\\hline
\multirow{2}{9.2cm}{以東人多益來告訴掃羅說：「大衛到了亞希米勒家。」那時，大衛作這訓誨詩，交於伶長。
1 勇士啊，你為何以作惡自誇？神的慈愛是常存的。
2 你的舌頭邪惡詭詐，好像剃頭刀，快利傷人。
3 你愛惡勝似愛善，又愛說謊，不愛說公義。（細拉）
4 詭詐的舌頭啊，你愛說一切毀滅的話！
5 神也要毀滅你，直到永遠。他要把你拿去，從你的帳篷中抽出，從活人之地將你拔出。（細拉）}
&義人要看見而懼怕，並要笑他，
7 說：「看哪！這就是那不以神為他力量的人，只倚仗他豐富的財物，在邪惡上堅立自己。」\\\cline{2-2}
&至於我，就像神殿中的青橄欖樹，我永永遠遠倚靠神的慈愛。
9 我要稱謝你，直到永遠，因為你行了這事。我也要在你聖民面前仰望你的名，這名本為美好。\\\hline
\end{tabular}
\end{table} 


\subsection{第53篇(大衛)：愚人不知有神吞吃神的百姓，神必救回被擄的子民 1-6}
愚頑人心裡說：「沒有神！」他們都是邪惡，行了可憎惡的罪孽，沒有一個人行善。
2 神從天上垂看世人，要看有明白的沒有，有尋求他的沒有。
3 他們各人都退後，一同變為汙穢；並沒有行善的，連一個也沒有。
4 作孽的沒有知識嗎？他們吞吃我的百姓如同吃飯一樣，並不求告神。
5 他們在無可懼怕之處，就大大害怕，因為神把那安營攻擊你之人的骨頭散開了。你使他們蒙羞，因為神棄絕了他們。
6 但願以色列的救恩從錫安而出！神救回他被擄的子民，那時雅各要快樂，以色列要歡喜。

\subsection{第54篇(大衛)：西弗人帮助掃羅尋索大衛，求主扶助 1-7}
\begin{table}[H]
\centering
\begin{tabular}{p{8.5cm}|p{3.500cm}|p{4.2500cm}}\hline\hline
\multicolumn{1}{c|}{外人攻擊，求神申冤}&\multicolumn{1}{c|}{求神滅絕仇敵}&\multicolumn{1}{c}{我要獻祭稱讚神}\\\hline
西弗人來對掃羅說：「大衛豈不是在我們那裡藏身嗎？」那時，大衛作這訓誨詩，交於伶長。用絲弦的樂器。
神啊，求你以你的名救我，憑你的大能為我申冤。
神啊，求你聽我的禱告，留心聽我口中的言語。
因為外人起來攻擊我，強暴人尋索我的命，他們眼中沒有神。（細拉）
&神是幫助我的，是扶持我命的。
5 他要報應我仇敵所行的惡，求你憑你的誠實滅絕他們。
&我要把甘心祭獻給你，耶和華啊，我要稱讚你的名，這名本為美好。
7 他從一切的急難中把我救出來，我的眼睛也看見了我仇敵遭報。\\\hline
\end{tabular}
\end{table} 

\subsection{第55篇(大衛)：受敵（密友）暴虐求其被滅施拯救 1-23}
\subsubsection{受密友暴虐求其被滅 1-15}
\begin{table}[H]
\centering
\begin{tabular}{p{3.5cm}|p{5.500cm}|p{7.7500cm}}\hline\hline
\multicolumn{1}{c|}{哀嘆原因}&\multicolumn{1}{c|}{我的心愿}&\multicolumn{1}{c}{密友的辱罵}\\\hline
神啊，求你留心聽我的禱告，不要隱藏不聽我的懇求。求你側耳聽我，應允我。我哀嘆不安，發聲唉哼，都因仇敵的聲音，惡人的欺壓。因為他們將罪孽加在我身上，發怒氣逼迫我。
&我心在我裡面甚是疼痛，死的驚惶臨到我身。恐懼戰兢歸到我身，驚恐漫過了我。我說：「但願我有翅膀像鴿子，我就飛去，得享安息。我必遠遊，宿在曠野。（細拉）
我必速速逃到避所，脫離狂風暴雨。」
主啊，求你吞滅他們，變亂他們的舌頭，因為我在城中見了強暴爭競的事。
&他們在城牆上晝夜繞行，在城內也有罪孽和奸惡。
邪惡在其中，欺壓和詭詐不離街市。
原來不是仇敵辱罵我，若是仇敵，還可忍耐。也不是恨我的人向我狂大，若是恨我的人，就必躲避他。
不料是你，你原與我平等，是我的同伴，是我知己的朋友！
我們素常彼此談論，以為甘甜，我們與群眾在神的殿中同行。
願死亡忽然臨到他們，願他們活活地下入陰間，因為他們的住處、他們的心中都是邪惡。\\\hline
\end{tabular}
\end{table} 

\subsubsection{继续求告因深信神必奖善罚惡施拯救 16-23}
\begin{table}[H]
\centering
\begin{tabular}{p{8.5cm}|p{9.2500cm}}\hline\hline
\multicolumn{1}{c|}{神待我}&\multicolumn{1}{c}{神待惡人}\\\hline
至於我，我要求告神，耶和華必拯救我。我要晚上、早晨、晌午哀聲悲嘆，他也必聽我的聲音。他救贖我命脫離攻擊我的人，使我得享平安，因為與我相爭的人甚多。&那沒有更變、不敬畏神的人，從太古常存的神必聽見而苦待他。他背了約，伸手攻擊與他和好的人。他的口如奶油光滑，他的心卻懷著爭戰；他的話比油柔和，其實是拔出來的刀。\\\hline
你要把你的重擔卸給耶和華，他必撫養你，他永不叫義人動搖。&神啊，你必使惡人下入滅亡的坑，流人血行詭詐的人必活不到半世，\\\hline
但我要倚靠你。&\\\hline
\end{tabular}
\end{table} 


\subsection{第56篇(大衛)： 哀訴仇敵之虐害，頌感神之救恩 1-13}
\subsubsection{非利士人在迦特拿住大衛，哀訴仇敵之虐害 1-9}
\begin{table}[H]
\centering
\begin{tabular}{p{11.75cm}|p{5.75cm}}\hline\hline
\multicolumn{1}{c|}{神待我}&\multicolumn{1}{c}{神待惡人}\\\hline
非利士人在迦特拿住大衛。那時，他作這金詩，交於伶長。調用遠方無聲鴿。
1 神啊，求你憐憫我，因為人要把我吞了，終日攻擊欺壓我。
2 我的仇敵終日要把我吞了，因逞驕傲攻擊我的人甚多。
3 我懼怕的時候要倚靠你。
4 我倚靠神，我要讚美他的話。我倚靠神，必不懼怕，血氣之輩能把我怎麼樣呢？
5 他們終日顛倒我的話，他們一切的心思都是要害我。
6 他們聚集、埋伏，窺探我的腳蹤，等候要害我的命。
&7 他們豈能因罪孽逃脫嗎？神啊，求你在怒中使眾民墮落！
8 我幾次流離，你都記數，求你把我眼淚裝在你的皮袋裡。這不都記在你冊子上嗎？
9 我呼求的日子，我的仇敵都要轉身退後。神幫助我，這是我所知道的。\\\hline
\end{tabular}
\end{table} 

\subsubsection{頌感神之救恩 10-13}
我倚靠神，我要讚美他的話！我倚靠耶和華，我要讚美他的話！
11 我倚靠神，必不懼怕，人能把我怎麼樣呢？
12 神啊，我向你所許的願在我身上，我要將感謝祭獻給你。
13 因為你救我的命脫離死亡，你豈不是救護我的腳不跌倒，使我在生命光中行在神面前嗎？

\subsection{第57篇(大衛藏洞裡逃避掃羅)： 被敵圍困靠主施援，要在萬民中稱謝神 1-11}
\begin{table}[H]
\centering
\begin{tabular}{p{8.7cm}|p{9cm}}\hline\hline
\multicolumn{1}{c|}{被敵圍困靠主施援1-4}&\multicolumn{1}{c}{敵人掉在自設的網羅，歌頌神的慈愛5-11}\\\hline
大衛逃避掃羅，藏在洞裡。那時，他作這金詩，交於伶長。調用休要毀壞。
神啊，求你憐憫我，憐憫我，因為我的心投靠你。我要投靠在你翅膀的蔭下，等到災害過去。
我要求告至高的神，就是為我成全諸事的神。
那要吞我的人辱罵我的時候，神從天上必施恩救我，也必向我發出慈愛和誠實。
我的性命在獅子中間，我躺臥在性如烈火的世人當中，他們的牙齒是槍、箭，他們的舌頭是快刀。
&神啊，願你崇高過於諸天，願你的榮耀高過全地！
他們為我的腳設下網羅，壓制我的心。他們在我面前挖了坑，自己反掉在其中。（細拉）
神啊，我心堅定，我心堅定！我要唱詩，我要歌頌！
我的靈(榮耀)啊，你當醒起！琴瑟啊，你們當醒起！我自己要極早醒起。
主啊，我要在萬民中稱謝你，在列邦中歌頌你。
因為你的慈愛高及諸天，你的誠實達到穹蒼。
神啊，願你崇高過於諸天，願你的榮耀高過全地！\\\hline
\end{tabular}
\end{table} 

\subsection{第58篇(大衛)： 求神降罰出母胎就行惡的世人 1-11}
\begin{table}[H]
\centering
\begin{tabular}{p{7.5cm}|p{9.8cm}}\hline\hline
\multicolumn{1}{c|}{出母胎就行惡的世人1-5}&\multicolumn{1}{c}{求神罰惡使義人歡喜6-11}\\\hline
世人哪，你們默然不語，真合公義嗎？施行審判，豈按正直嗎？
不然！你們是心中作惡，你們在地上稱出你們手所行的強暴。
惡人一出母胎，就與神疏遠；一離母腹，便走錯路，說謊話。
他們的毒氣好像蛇的毒氣，他們好像塞耳的聾虺，
不聽行法術的聲音，雖用極靈的咒語，也是不聽。
&神啊，求你敲碎他們口中的牙！耶和華啊，求你敲掉少壯獅子的大牙！
願他們消滅如急流的水一般，他們瞅準射箭的時候，願箭頭彷彿砍斷。
願他們像蝸牛消化過去，又像婦人墜落未見天日的胎。
你們用荊棘燒火，鍋還未熱，他要用旋風把青的和燒著的一齊颳去。
義人見仇敵遭報就歡喜，要在惡人的血中洗腳。
 因此，人必說：「義人誠然有善報，在地上果有施行判斷的神。」\\\hline
\end{tabular}
\end{table} 

\subsection{第59篇(掃羅窺探大衛房屋要殺他)： 求主救援脫離敵害，并歌頌神的救恩1-17}
\subsubsection{大衛因掃羅差凶手窺探房屋,求主救援脫離敵害 1-5}
掃羅打發人窺探大衛的房屋，要殺他。那時，大衛作這金詩，交於伶長。調用休要毀壞。
1 我的神啊，求你救我脫離仇敵，把我安置在高處，得脫那些起來攻擊我的人。
2 求你救我脫離作孽的人，和喜愛流人血的人。

3 因為他們埋伏要害我的命，有能力的人聚集來攻擊我。耶和華啊，這不是為我的過犯，也不是為我的罪愆。
4 我雖然無過，他們預備整齊，跑來攻擊我。

求你興起，鑒察、幫助我！
5 萬軍之神耶和華，以色列的神啊，求你興起，懲治萬邦，不要憐憫行詭詐的惡人。（細拉）

\subsubsection{相信必見仇敵遭报，并歌頌神的救恩 6-17}
\begin{table}[H]
\centering
\begin{tabular}{p{10.5cm}|p{6.8cm}}\hline\hline
\multicolumn{1}{c|}{惡人所行和求神施报}&\multicolumn{1}{c}{必見仇敵遭报并歌頌神的救恩}\\\hline
他們晚上轉回，叫號如狗，圍城繞行。
他們口中噴吐惡言，嘴裡有刀，他們說：「有誰聽見？」
但你耶和華必笑話他們，你要嗤笑萬邦。
&我的力量啊，我必仰望你，因為神是我的高臺。
我的神要以慈愛迎接我，神要叫我看見我仇敵遭報。\\\hline
不要殺他們，恐怕我的民忘記。主啊，你是我們的盾牌，求你用你的能力使他們四散，且降為卑。
因他們口中的罪和嘴裡的言語，並咒罵虛謊的話，願他們在驕傲之中被纏住了。
&\multirow{2}{6.9cm}{但我要歌頌你的力量，早晨要高唱你的慈愛，因為你做過我的高臺，在我急難的日子做過我的避難所。
我的力量啊，我要歌頌你！因為神是我的高臺，是賜恩於我的神。}\\\cline{1-1}
求你發怒，使他們消滅，以至歸於無有，叫他們知道神在雅各中間掌權，直到地極。到了晚上，任憑他們轉回，任憑他們叫號如狗，圍城繞行。
他們必走來走去，尋找食物，若不得飽，就終夜在外。&\\\hline
\end{tabular}
\end{table} 


\subsection{第60篇(大衛)：被神丟棄而遭破敗，求神幫助踐踏敵人 1-12}
\subsubsection{因被神丟棄而遭破敗向神哀訴 1-4}
大衛與兩河間的亞蘭並瑣巴的亞蘭爭戰的時候，約押轉回，在鹽谷攻擊以東，殺了一萬二千人。那時，大衛作這金詩叫人學習，交於伶長。調用為證的百合花。
1 神啊，你丟棄了我們，使我們破敗，你向我們發怒，求你使我們復興。
2 你使地震動，而且崩裂，求你將裂口醫好，因為地搖動。
3 你叫你的民遇見艱難，你叫我們喝那使人東倒西歪的酒。
4 你把旌旗賜給敬畏你的人，可以為真理揚起來。（細拉）
\subsubsection{倚靠神踐踏敵人展大能 5-12}
求你應允我們，用右手拯救我們，好叫你所親愛的人得救。
6 神已經指著他的聖潔說(應許我)：「我要歡樂，我要分開示劍，丈量疏割谷。
7 基列是我的，瑪拿西也是我的。以法蓮是護衛我頭的，猶大是我的杖。
8 摩押是我的沐浴盆，我要向以東拋鞋。非利士啊，你還能因我歡呼嗎？」
9 誰能領我進堅固城？誰能引我到以東地？
10 神啊，你不是丟棄了我們嗎？神啊，你不和我們的軍兵同去嗎？
11 求你幫助我們攻擊敵人，因為人的幫助是枉然的。
12 我們倚靠神才得施展大能，因為踐踏我們敵人的就是他。

\subsection{第61篇(大衛)： 投靠神必蒙恩佑到永遠 1-8}
神啊，求你聽我的呼求，側耳聽我的禱告。
2 我心裡發昏的時候，我要從地極求告你，求你領我到那比我更高的磐石。
3 因為你做過我的避難所，做過我的堅固臺，脫離仇敵。
4 我要永遠住在你的帳幕裡，我要投靠在你翅膀下的隱密處。（細拉）

5 神啊，你原是聽了我所許的願；你將產業賜給敬畏你名的人。
6 你要加添王的壽數，他的年歲必存到世世。
7 他必永遠坐在神面前，願你預備慈愛和誠實保佑他。
8 這樣，我要歌頌你的名直到永遠，好天天還我所許的願。

\subsection{第62篇(大衛)： 不畏人欺默候神，能力慈愛屬乎神 1-12}
我的心默默無聲，專等候神，我的救恩是從他而來。
2 唯獨他是我的磐石、我的拯救，他是我的高臺，我必不很動搖。

3 你們大家攻擊一人，把他毀壞，如同毀壞歪斜的牆、將倒的壁，要到幾時呢？
4 他們彼此商議，專要從他的尊位上把他推下。他們喜愛謊話，口雖祝福，心卻咒詛。（細拉）

5 我的心哪，你當默默無聲，專等候神，因為我的盼望是從他而來。
6 唯獨他是我的磐石、我的拯救，他是我的高臺，我必不動搖。
7 我的拯救，我的榮耀，都在乎神；我力量的磐石，我的避難所，都在乎神。
8 你們眾民當時時倚靠他，在他面前傾心吐意，神是我們的避難所。（細拉）

9 下流人真是虛空，上流人也是虛假，放在天平裡就必浮起，他們一共比空氣還輕。
10 不要仗勢欺人，也不要因搶奪而驕傲，若財寶加增，不要放在心上。

11 神說了一次兩次，我都聽見，就是能力都屬乎神。
12 主啊，慈愛也是屬乎你，因為你照著各人所行的報應他。

\subsection{第63篇(大衛)：尋求神者必得飽足，殺人的必被刀劍所滅  1-11}

大衛在猶大曠野的時候，作了這詩。
1 神啊，你是我的神，我要切切地尋求你！在乾旱疲乏無水之地，我渴想你，我的心切慕你。
2 我在聖所中曾如此瞻仰你，為要見你的能力和你的榮耀。
3 因你的慈愛比生命更好，我的嘴唇要頌讚你。

4 我還活的時候要這樣稱頌你，我要奉你的名舉手。
5 我在床上記念你，在夜更的時候思想你，我的心就像飽足了骨髓肥油，
6 我也要以歡樂的嘴唇讚美你。
7 因為你曾幫助我，我就在你翅膀的蔭下歡呼。
8 我心緊緊地跟隨你，你的右手扶持我。

9 但那些尋索要滅我命的人，必往地底下去。
10 他們必被刀劍所殺，被野狗所吃。
11 但是王必因神歡喜，凡指著他發誓的，必要誇口，因為說謊之人的口必被塞住。


\subsection{第64篇(大衛)： 求主保護脫離仇敵，使義人歡喜誇口 1-10}
\begin{table}[H]
\centering
\begin{tabular}{p{3.7cm}|p{6.5cm}|p{6.4cm}}\hline\hline
\multicolumn{1}{c|}{求神保護性命}&\multicolumn{1}{c|}{惡人所行}&\multicolumn{1}{c}{惡人、眾人，義人}\\\hline
\multirow{3}{3.8cm}{神啊，我哀嘆的時候，求你聽我的聲音。求你保護我的性命，不受仇敵的驚恐。
求你把我隱藏，使我脫離作惡之人的暗謀和作孽之人的擾亂。}
&\multirow{3}{6.6cm}{他們磨舌如刀，發出苦毒的言語，好像比準了的箭，
4 要\textcolor{red}{在暗地射}完全人。他們忽然射他，並不懼怕。
5 他們彼此勉勵設下惡計，他們商量暗設網羅，說：「誰能看見？」
6 他們圖謀奸惡，說：「我們是極力圖謀的！」他們各人的意念心思是深的。}
&但\textcolor{red}{神要射}他們，他們忽然被箭射傷。
他們必然絆跌，被自己的舌頭所害，\\\cline{3-3}
&&凡看見他們的必都搖頭。眾人都要害怕，要傳揚神的工作，並且明白他的作為。\\\cline{3-3}
&&義人必因耶和華歡喜，並要投靠他。凡心裡正直的人，都要誇口。\\\hline
\end{tabular}
\end{table} 



\subsection{第65篇(大衛)：讚主赦罪之恩,是萬民的倚靠,恩澤萬物 1-13}
\subsubsection{選民——頌讚主赦罪之恩,所揀選的人得親近主 1-4}
神啊，錫安的人都等候讚美你，所許的願也要向你償還。
2 聽禱告的主啊，凡有血氣的都要來就你。
3 罪孽勝了我，至於我們的過犯，你都要赦免。
4 你所揀選使他親近你，住在你院中的，這人便為\textcolor{red}{有福}！我們必因你居所，你聖殿的美福\textcolor{red}{知足}了。
\subsubsection{萬民——平靜萬民的喧嘩,使地極的人都歡呼 5-8}
5 拯救我們的神啊，你必以威嚴秉公義應允我們，你本是一切地極和海上遠處的人所倚靠的。
6 他既以大能束腰，就用力量安定諸山，
7 使諸海的響聲和其中波浪的響聲，並萬民的喧嘩，都平靜了。
8 住在地極的人因你的神蹟懼怕，你使日出日落之地都歡呼。
\subsubsection{自然——述神之恩澤及萬物 9-13}
你眷顧地，降下透雨，使地大得肥美；神的河滿了水。你這樣澆灌了地，好為人預備五穀。
10 你澆透地的犁溝，潤平犁脊，降甘霖使地軟和，其中發長的蒙你賜福。
11 你以恩典為年歲的冠冕，你的路徑都滴下脂油。
12 滴在曠野的草場上，小山以歡樂束腰，
13 草場以羊群為衣，谷中也長滿了五穀，這一切都歡呼歌唱。

\subsection{第66篇：頌讚神之大能,進神殿獻祭還願,稱頌神聽允其祈 1-20}
\subsubsection{頌讚神之大能 1-12}
\begin{table}[H]
\centering
\begin{tabular}{p{7.05cm}|p{10.5cm}}\hline\hline
\multicolumn{1}{c|}{歌頌讚美神名的榮耀}&\multicolumn{1}{c}{细数神所行的}\\\hline
全地都當向神歡呼！
歌頌他名的榮耀，用讚美的言語將他的榮耀發明！
當對神說：「你的作為何等可畏！因你的大能，仇敵要投降你。
全地要敬拜你，歌頌你，要歌頌你的名。」（細拉）
&你們來看神所行的，他向世人所做之事是可畏的。
他將海變成乾地，眾民步行過河，我們在那裡因他歡喜。
他用權能治理萬民直到永遠，他的眼睛鑒察列邦，悖逆的人不可自高。（細拉）\\\hline
萬民哪，你們當稱頌我們的神，使人得聽讚美他的聲音。
&他使我們的性命存活，也不叫我們的腳搖動。
神啊，你曾試驗我們，熬煉我們如熬煉銀子一樣。
你使我們進入網羅，把重擔放在我們的身上。
你使人坐車軋我們的頭，我們經過水火，你卻使我們到豐富之地。\\\hline
\end{tabular}
\end{table} 

\subsubsection{進神殿獻祭還願,稱頌神聽允其祈 13-20}
\begin{table}[H]
\centering
\begin{tabular}{p{6.25cm}|p{11.cm}}\hline\hline
\multicolumn{1}{c|}{燔祭進神殿還願13-15}&\multicolumn{1}{c}{稱頌神聽允其祈16-20}\\\hline
\multirow{2}{6.3cm}{我要用燔祭進你的殿，向你還我的願，
就是在急難時我嘴唇所發的，口中所許的。
我要把肥牛做燔祭，將公羊的香祭獻給你，又把公牛和山羊獻上。（細拉）}
&凡敬畏神的人，你們都來聽，我要述說他為我所行的事。
我曾用口求告他，我的舌頭也稱他為高。
我若心裡注重罪孽，主必不聽。\\\cline{2-2}
&但神實在聽見了，他側耳聽了我禱告的聲音。
神是應當稱頌的！他並沒有推卻我的禱告，也沒有叫他的慈愛離開我。\\\hline
\end{tabular}
\end{table} 

\subsection{第67篇 勸萬民稱讚神的救恩,公正審判和賜福 1-7}
願神憐憫我們，賜福於我們，用臉光照我們，（細拉）
2 好叫世界得知你的道路，萬國得知你的\textcolor{red}{救恩}。
3 神啊，願列邦稱讚你，願萬民都稱讚你！
4 願萬國都快樂歡呼，因為你必按\textcolor{red}{公正審判}萬民，引導世上的萬國。（細拉）
5 神啊，願列邦稱讚你，願萬民都稱讚你！
6 地已經出了土產，神，就是我們的神，要\textcolor{red}{賜福}於我們。
7 神要賜福於我們，地的四極都要敬畏他。

\subsection{第68篇(基督国度的大卫之约)：讚美神救赎其為奴和被掳之民 1-35}
\subsubsection{求主興起使惡人消滅義人歡欣 1-6}
願神興起，使他的仇敵四散，叫那恨他的人從他面前逃跑。
2 他們被驅逐，如煙被風吹散；惡人見神之面而消滅，如蠟被火熔化。

3 唯有義人必然歡喜，在神面前高興快樂。
4 你們當向神唱詩，歌頌他的名，為那坐車行過曠野的修平大路——他的名是耶和華，要在他面前歡樂！

5 神在他的聖所做孤兒的父，做寡婦的申冤者。
6 神叫孤獨的有家，使被囚的出來享福，唯有悖逆的住在乾燥之地。

\subsubsection{讚美神眷顧其民（基督的救赎） 7-18}
神啊，你曾在你百姓前頭出來，在曠野行走。（細拉）
8 那時地見神的面而震動，天也落雨，西奈山見以色列神的面也震動。
9 神啊，你降下大雨，你產業以色列疲乏的時候，你使他堅固。
10 你的會眾住在其中，神啊，你的恩惠是為困苦人預備的。
11 主發命令，傳好信息的婦女成了大群。
12 統兵的君王逃跑了，逃跑了！在家等候的婦女，分受所奪的。
13 你們安臥在羊圈的時候，好像鴿子的翅膀鍍白銀，翎毛鍍黃金一般。
14 全能者在境內趕散列王的時候，勢如飄雪在撒們。
15 巴珊山是神的山，巴珊山是多峰多嶺的山。
16 你們多峰多嶺的山哪，為何斜看神所願居住的山？耶和華必住這山，直到永遠。
17 神的車輦累萬盈千，主在其中，好像在西奈聖山一樣。
18 你已經升上高天，擄掠仇敵。你在人間，就是在悖逆的人間，受了供獻，叫耶和華神可以與他們同住。
\subsubsection{稱頌神神使被掳的归回，列國要向神歌唱 19-35}
天天背負我們重擔的主，就是拯救我們的神，是應當稱頌的！（細拉）
20 神是為我們施行諸般救恩的神，人能脫離死亡是在乎主耶和華。
21 但神要打破他仇敵的頭，就是那常犯罪之人的髮頂。
22 主說：「我要使眾民從巴珊而歸，使他們從深海而回，
23 使你打碎仇敵，你的腳踹在血中，使你狗的舌頭從其中得份。」
24 神啊，你是我的神，我的王，人已經看見你行走，進入聖所。
25 歌唱的行在前，作樂的隨在後，都在擊鼓的童女中間。
26 從以色列源頭而來的，當在各會中稱頌主神。
27 在那裡有統管他們的小便雅憫，有猶大的首領和他們的群眾，有西布倫的首領，有拿弗他利的首領。
28 以色列的能力是神所賜的，神啊，求你堅固你為我們所成全的事。
29 因你耶路撒冷的殿，列王必帶貢物獻給你。
30 求你叱喝蘆葦中的野獸和群公牛，並列邦中的牛犢，把銀塊踹在腳下；神已經趕散好爭戰的列邦。

31 埃及的公侯要出來朝見神，古實人要急忙舉手禱告。
32 世上的列國啊，你們要向神歌唱！願你們歌頌主，
33 歌頌那自古駕行在諸天以上的主！他發出聲音，是極大的聲音。
34 你們要將能力歸給神，他的威榮在以色列之上，他的能力是在穹蒼。
35 神啊，你從聖所顯為可畏。以色列的神是那將力量、權能賜給他百姓的。神是應當稱頌的！

\subsection{第69篇(大衛:基督的受苦-克西马尼):哀陳仇敵強盛求神懲罰，願神要拯救錫安使其民讚美神名 1-36}
\subsubsection{哀陳苦況求主眷顧 1-21}
神啊，求你救我！因為眾水要淹沒我。
2 我陷在深淤泥中，沒有立腳之地；我到了深水中，大水漫過我身。
3 我因呼求困乏，喉嚨發乾；我因等候神，眼睛失明。
4 無故恨我的，比我頭髮還多；無理與我為仇，要把我剪除的，甚為強盛；我沒有搶奪的，要叫我償還。

5 神啊，我的愚昧你原知道，我的罪愆不能隱瞞。
6 萬軍的主耶和華啊，求你叫那等候你的不要因我蒙羞。以色列的神啊，求你叫那尋求你的不要因我受辱。

7 因我為你的緣故受了辱罵，滿面羞愧。
8 我的弟兄看我為外路人，我的同胞看我為外邦人。

9 因我為你的殿心裡焦急，如同火燒，並且辱罵你人的辱罵都落在我身上。
10 我哭泣，以禁食刻苦我心，這倒算為我的羞辱。
11 我拿麻布當衣裳，就成了他們的笑談。
12 坐在城門口的談論我，酒徒也以我為歌曲。


13 但我在悅納的時候，向你耶和華祈禱。神啊，求你按你豐盛的慈愛，憑你拯救的誠實，應允我！
14 求你搭救我出離淤泥，不叫我陷在其中；求你使我脫離那些恨我的人，使我出離深水。
15 求你不容大水漫過我，不容深淵吞滅我，不容坑坎在我以上合口。
16 耶和華啊，求你應允我，因為你的慈愛本為美好。求你按你豐盛的慈悲，回轉眷顧我。
不要掩面不顧你的僕人，我是在急難之中，求你速速地應允我。
18 求你親近我，救贖我，求你因我的仇敵把我贖回。
19 你知道我受的辱罵、欺凌、羞辱，我的敵人都在你面前。
20 辱罵傷破了我的心，我又滿了憂愁。我指望有人體恤，卻沒有一個；我指望有人安慰，卻找不著一個。
21 他們拿苦膽給我當食物，我渴了，他們拿醋給我喝。
\subsubsection{求神懲罰仇敵 22-28}
願他們的筵席在他們面前變為網羅，在他們平安的時候變為機檻。
23 願他們的眼睛昏矇，不得看見，願你使他們的腰常常戰抖。
24 求你將你的惱恨倒在他們身上，叫你的烈怒追上他們。
25 願他們的住處變為荒場，願他們的帳篷無人居住。
26 因為你所擊打的，他們就逼迫；你所擊傷的，他們戲說他的愁苦。
27 願你在他們的罪上加罪，不容他們在你面前稱義。
28 願他們從生命冊上被塗抹，不得記錄在義人之中。
\subsubsection{願神要拯救錫安，他的民讚美神的名 29-36}
29 但我是困苦憂傷的，神啊，願你的救恩將我安置在高處！
30 我要以詩歌讚美神的名，以感謝稱他為大。
31 這便叫耶和華喜悅，勝似獻牛，或是獻有角有蹄的公牛。
32 謙卑的人看見了，就喜樂。尋求神的人，願你們的心甦醒。
33 因為耶和華聽了窮乏人，不藐視被囚的人。
34 願天和地，洋海和其中一切的動物，都讚美他。
35 因為神要拯救錫安，建造猶大的城邑，他的民要在那裡居住，得以為業。
36 他僕人的後裔要承受為業，愛他名的人也要住在其中。


\subsection{第70篇(大衛)： 求耶和華速速救助使敵蒙羞 1-5}
神啊，求你快快搭救我！耶和華啊，求你速速幫助我！
2 願那些尋索我命的抱愧蒙羞，願那些喜悅我遭害的退後受辱。
3 願那些對我說「啊哈！啊哈！」的，因羞愧退後。
4 願一切尋求你的因你高興歡喜，願那些喜愛你救恩的常說：「當尊神為大！」
5 但我是困苦窮乏的，神啊，求你速速到我這裡來！你是幫助我的，搭救我的，耶和華啊，求你不要耽延！

\subsection{第71篇：年幼到年老要靠主扶助，讚美其公義救恩大能奇行 1-24}
\subsubsection{自幼倚主求其扶助 1-8}
耶和華啊，我投靠你，求你叫我永不羞愧。
2 求你憑你的公義搭救我，救拔我；側耳聽我，拯救我。
3 求你做我常住的磐石，你已經命定要救我，因為你是我的巖石，我的山寨。
4 我的神啊，求你救我脫離惡人的手，脫離不義和殘暴之人的手。
5 主耶和華啊，你是我所盼望的，從我年幼你是我所倚靠的。
6 我從出母胎被你扶持，使我出母腹的是你，我必常常讚美你。
7 許多人以我為怪，但你是我堅固的避難所。
8 你的讚美，你的榮耀，終日必滿了我的口。
\subsubsection{年老時求主不要丟棄,並要越發讚美 9-16}
9 我年老的時候，求你不要丟棄我；我力氣衰弱的時候，求你不要離棄我。
10 我的仇敵議論我，那些窺探要害我命的彼此商議，
11 說：「神已經離棄他，我們追趕他，捉拿他吧！因為沒有人搭救。」
12 神啊，求你不要遠離我！我的神啊，求你速速幫助我！
13 願那與我性命為敵的羞愧被滅，願那謀害我的受辱蒙羞。
我卻要常常盼望，並要越發讚美你。
15 我的口終日要述說你的公義和你的救恩，因我不計其數。
16 我要來說主耶和華大能的事，我單要提說你的公義。

\subsubsection{終生要述其公義救恩大能奇行 17-24}
17 神啊，自我年幼時你就教訓我，直到如今我傳揚你奇妙的作為。
18 神啊，我到年老髮白的時候，求你不要離棄我，等我將你的能力指示下代，將你的大能指示後世的人。
19 神啊，你的公義甚高；行過大事的神啊，誰能像你？
20 你是叫我們多經歷重大急難的，必使我們復活，從地的深處救上來。
21 求你使我越發昌大，又轉來安慰我。
22 我的神啊，我要鼓瑟稱讚你，稱讚你的誠實！以色列的聖者啊，我要彈琴歌頌你！
23 我歌頌你的時候，我的嘴唇和你所贖我的靈魂都必歡呼。
24 並且我的舌頭必終日講論你的公義，因為那些謀害我的人已經蒙羞受辱了。

\subsection{第72篇(所羅門)：王按公義審判掌權,被万国稱頌 1-20}
\begin{table}[H]
\centering
%\begin{tabular}{p{3.cm}|p{2.9cm}|p{2cm}|p{3.5cm}|p{4.75cm}}\hline\hline
\begin{tabular}{p{8.5em}|p{8.2em}|p{5.67em}|p{10.3em}|p{13.6em}}\hline\hline
\multicolumn{1}{c|}{按公義審判}&\multicolumn{1}{c|}{執掌權柄}&\multicolumn{1}{c|}{萬國侍奉他}&\multicolumn{1}{c|}{憐恤救贖}&\multicolumn{1}{c}{萬國稱頌他的名}\\\hline
神啊，求你將判斷的權柄賜給王，將公義賜給王的兒子。
他要按公義審判你的民，按公平審判你的困苦人。
大山小山，都要因公義使民得享平安。
他必為民中的困苦人申冤，拯救窮乏之輩，壓碎那欺壓人的。
&太陽還存，月亮還在，人要敬畏你，直到萬代。
他必降臨，像雨降在已割的草地上，如甘霖滋潤田地。
在他的日子義人要發旺，大有平安，好像月亮長存。
他要執掌權柄，從這海直到那海，從大河直到地極。
&住在曠野的必在他面前下拜，他的仇敵必要舔土。
他施和海島的王要進貢，示巴和西巴的王要獻禮物。
諸王都要叩拜他，萬國都要侍奉他。
&因為窮乏人呼求的時候，他要搭救；沒有人幫助的困苦人，他也要搭救。
他要憐恤貧寒和窮乏的人，拯救窮苦人的性命。
他要救贖他們脫離欺壓和強暴，他們的血在他眼中看為寶貴，
他們要存活。示巴的金子要奉給他，人要常常為他禱告，終日稱頌他。
&在地的山頂上五穀必然茂盛 (有一把五穀)，所結的穀實要響動如黎巴嫩的樹林，城裡的人要發旺如地上的草。
他的名要存到永遠，要留傳如日之久。人要因他蒙福，萬國要稱他有福。
獨行奇事的耶和華以色列的神是應當稱頌的！
他榮耀的名也當稱頌，直到永遠！願他的榮耀充滿全地！阿們，阿們。
耶西的兒子大衛的祈禱完畢。\\\hline
\end{tabular}
\end{table} 


